\hypersetup{pageanchor=false}
\begin{titlepage}
\centering
\vspace*{-0.4cm}

\begin{tcolorbox}[colback=Primary,colframe=Primary,boxrule=0pt,arc=2.5mm,
  width=\textwidth,left=2.5mm,right=2.5mm,top=2.5mm,bottom=2.5mm]
\centering{\color{white}\bfseries\large UNIVERSIDAD DE LAS AMÉRICAS PUEBLA}
\end{tcolorbox}

\vspace{0.1cm}
\IfFileExists{images/logo.png}{\includegraphics[width=0.13\textwidth]{images/logo.png}}{}
\vspace{0.1cm}

\begin{tcolorbox}[colback=SoftGray,colframe=Primary,boxrule=1pt,arc=2.5mm,
  width=\textwidth,left=3mm,right=3mm,top=2.5mm,bottom=2.5mm]
\centering
{\Large\bfseries\color{Secondary} SISTEMAS EMBEBIDOS}\par
\vspace{0.08cm}
{\normalsize\color{TextDark}\clave\ \,$\cdot$\,\ Reporte del proyecto final y práctica 5}\par
\vspace{0.15cm}
\begin{tcolorbox}[colback=white,colframe=Secondary,boxrule=1pt,arc=2mm,
  width=0.95\textwidth,left=2.4mm,right=2.4mm,top=2mm,bottom=2mm]
\centering{\LARGE\bfseries\color{Primary}\titulo\\[4pt]
\normalsize\mdseries\color{TextDark}\subtitulo}
\end{tcolorbox}
\end{tcolorbox}

\vspace{0.15cm}

\begin{minipage}[t]{0.50\textwidth}
\begin{tcolorbox}[colback=SoftBlue,colframe=Secondary,boxrule=1pt,arc=2mm,
  title=\textbf{Integrantes}]
\small
\setlength{\tabcolsep}{2.5pt}
\begin{tabular}{@{}l l@{}}
\autorA & ID: \idA \\
\autorB & ID: \idB \\
\autorC & ID: \idC \\[2pt]
\multicolumn{2}{@{}l}{\textbf{Fecha:} \fechadoc}
\end{tabular}
\end{tcolorbox}
\end{minipage}
\hfill
\begin{minipage}[t]{0.47\textwidth}
\begin{tcolorbox}[colback=AccentSoft,colframe=Accent,boxrule=1pt,arc=2mm,
  title=\textbf{Código, proyectos y evidencias}]
\small\raggedright
Repositorio: \url{\repositorio}\par
\vspace{2pt}
Verilog, bancos de prueba, proyectos de Quartus, layout en SKY130 y los programas que
generan cada figura de este documento.
\end{tcolorbox}
\end{minipage}

\vfill
\begin{tcolorbox}[colback=white,colframe=Primary,boxrule=0.9pt,arc=2mm,
  width=\textwidth,title=\textbf{Resumen},colbacktitle=Primary,coltitle=white,
  left=2mm,right=2mm,top=1mm,bottom=1mm]
\footnotesize\justifying
Después de un sismo hay personas que quedan en huecos de unos centímetros, a los que no llega un perro, una cámara ni un
robot. En este proyecto diseñamos, verificamos y prototipamos en FPGA el sistema digital de una mochila para que la
cucaracha silbadora de Madagascar (\especie) ayude a buscarlas, junto con la red con la que varias mochilas le avisan a un
operador. A partir de datos publicados fijamos una carga útil de 2 a 3~g y un estímulo que no requiere cirugía y que
limitamos con un presupuesto para no habituar al insecto. El núcleo de la mochila cuenta los píxeles calientes de una
matriz térmica y decide con una máquina de ocho estados. Ocupa 477 elementos lógicos en la MAX~10 de la DE10-Lite y dio lo
mismo que su modelo de referencia en los 11\,424 cuadros que probamos. Con él armamos una red de tres tarjetas, dos
mochilas y un receptor, unidas por UART con tramas de 16 bytes, CRC-16, número de secuencia y retransmisión. Cada nodo
emite un tono, reconoce a sus vecinos con filtros de Goertzel y calcula distancias exactas con una raíz cuadrada entera; el
receptor triangula a la persona con un CORDIC de 16 iteraciones y se equivoca 6.2~cm a 31.95~m. La red pasó los nueve
escenarios de prueba, entre ellos el ruido en la línea y el corte del alcance en 32~m, y ocupa 3\,933 elementos lógicos por
mochila a 55~MHz. También llevamos el núcleo a un layout en SKY130 con el flujo abierto LibreLane: mide 0.050~mm$^2$, no
tiene errores de DRC ni de LVS, alcanza unos 72~MHz en la esquina lenta y gasta 0.58~mA a 1.8~V. En una simulación de
\numMisiones{} misiones, diez insectos guiados encontraron a la persona en \pGuiadaSesenta~\% de las misiones de una hora,
contra \pPasivaSesenta~\% sin estímulos. El sensor y el canal acústico están simulados dentro de la FPGA, no hicimos pruebas
con insectos y la ley mexicana pide mantener a esta especie en confinamiento.

\vspace{2pt}
\textbf{Palabras clave:} FPGA, insecto cíborg, búsqueda y rescate, CORDIC, filtro de Goertzel,
CRC-16, triangulación, SKY130.
\end{tcolorbox}
\end{titlepage}
